% !TeX root = main.tex
\changecolours{ScoutNavy}
\section[CRUD do Zero: Compass e mongosh]{Aprendendo do Zero: CRUD Passo a Passo no Compass e \texttt{mongosh}}

% ------------------------------------------------------------------------------
\begin{frame}[fragile]{Passo 1: Criando o Banco e Consultando Coleções}
  \vspace{0.05cm}
  \begin{columns}[t]
    \begin{column}{0.48\textwidth}
      \alert{1. Pela Interface Gráfica (MongoDB Compass)}\par
      \itemseps[topsepi=0.2mm,itemsepi=0.5mm,labelsep=1.5mm,leftmargini=2.5mm]
      \scriptsize
      \begin{enumerate}
        \item Conecte em \texttt{mongodb://localhost:27017}.
        \item No menu esquerdo (\textit{Databases}), clique em \textbf{\texttt{+} (Create database)}.
        \item Preencha \textbf{Database Name:} \texttt{mongodb\_books} e \textbf{Collection Name:} \texttt{books}.
        \item \textbf{Consulte a lista de coleções:} Clique sobre \texttt{mongodb\_books} no menu lateral e verifique a coleção \texttt{books} na aba \textbf{Collections}.
      \end{enumerate}
    \end{column}
    \begin{column}{0.48\textwidth}
      \alert{2. Pelo Terminal Interativo (\texttt{mongosh})}\par
      \begin{lstlisting}[language=mongosh]
// Listar todos os bancos existentes:
show dbs

// Selecionar/criar o banco mongodb_books:
use mongodb_books

// Criar explicitamente a colecao books:
db.createCollection("books")

// Consultar a lista de colecoes do banco:
show collections
// Equivalente em JavaScript:
db.getCollectionNames()
      \end{lstlisting}
    \end{column}
  \end{columns}
\end{frame}

% ------------------------------------------------------------------------------
\begin{frame}[fragile]{Passo 2: Inserindo (\textit{Create}) e Consultando (\textit{Read}) o 1.º Documento}
  \vspace{0.05cm}
  \begin{columns}[t]
    \begin{column}{0.48\textwidth}
      \alert{1. Inserção Manual no MongoDB Compass}\par
      \scriptsize
      Em \texttt{books}, clique em \textbf{Add Data $\to$ Insert Document}, cole o JSON e clique em \textbf{Insert}:
      \begin{lstlisting}[language=mongosh]
{
  "isbn": "101",
  "name": "MongoDB 7.0 Guia Pratico",
  "price": 30
}
      \end{lstlisting}
      \scriptsize
      $\bullet$ \textbf{Consultar a lista de documentos:} Na aba \textbf{Documents}, veja que o MongoDB gerou o campo \texttt{\_id} (\texttt{ObjectId}) automaticamente! Alterne entre as visões \textbf{List}, \textbf{JSON} e \textbf{Table}.
    \end{column}
    \begin{column}{0.48\textwidth}
      \alert{2. Equivalente no \texttt{mongosh} (\texttt{insertOne} e \texttt{find})}\par
      \begin{lstlisting}[language=mongosh]
// Inserir 1 documento:
db.books.insertOne({
  isbn: "101",
  name: "MongoDB 7.0 Guia Pratico",
  price: 30
})
// Retorno:
// { acknowledged: true,
//   insertedId: ObjectId("671ec880...") }

// Consultar todos os documentos da colecao:
db.books.find()
      \end{lstlisting}
    \end{column}
  \end{columns}
\end{frame}

% ------------------------------------------------------------------------------
\begin{frame}[fragile]{Passo 3: Subdocumentos Aninhados e Busca com \textit{Dot Notation}}
  \vspace{0.05cm}
  \begin{columns}[t]
    \begin{column}{0.48\textwidth}
      \alert{1. Inserindo Documento com Objeto \texttt{meta}}\par
      \scriptsize
      Em \textbf{Add Data $\to$ Insert Document} no Compass (após inserir, consulte a lista de documentos e expanda \texttt{meta}):
      \begin{lstlisting}[language=mongosh]
{
  "isbn": "201",
  "name": "Modelagem de Documentos NoSQL",
  "meta": {
    "version": "MongoDB 6.0"
  },
  "price": 50
}
      \end{lstlisting}
    \end{column}
    \begin{column}{0.48\textwidth}
      \alert{2. Filtrando no Compass e no \texttt{mongosh}}\par
      \scriptsize
      Na barra \textbf{Filter} do Compass (ou dentro de \texttt{db.books.find(...)}):
      \begin{lstlisting}[language=mongosh]
// 1. Busca exata por atributo de topo:
{ "isbn": "101" }

// 2. Dot Notation em subdocumento:
{ "meta.version": "MongoDB 6.0" }

// 3. Operador de comparacao ($gt = maior que):
{ "price": { "$gt": 35 } }

// 4. Expressao Regular case-insensitive:
{ "name": { "$regex": "mongo", "$options": "i" } }
      \end{lstlisting}
    \end{column}
  \end{columns}
\end{frame}

% ------------------------------------------------------------------------------
\begin{frame}[fragile]{Passo 4: Atualizando (\textit{Update}) e Removendo (\textit{Delete}) Documentos}
  \vspace{0.05cm}
  \begin{columns}[t]
    \begin{column}{0.48\textwidth}
      \alert{1. Atualizando com \texttt{updateOne} e \texttt{\$set}}\par
      \begin{lstlisting}[language=mongosh]
// Atualizar o preco do livro de ISBN "101":
db.books.updateOne(
  { "isbn": "101" },          // Filtro de busca
  { "$set": { "price": 45 } } // Operador atomico
)
// Retorno:
// { acknowledged: true,
//   matchedCount: 1, modifiedCount: 1 }
      \end{lstlisting}
      \scriptsize
      $\bullet$ \textbf{No Compass:} Passe o mouse sobre o documento, clique no lápis (\textbf{Edit Document}), altere \texttt{price} para \texttt{45}, clique em \textbf{Update} e consulte a lista de documentos!
    \end{column}
    \begin{column}{0.48\textwidth}
      \alert{2. Removendo com \texttt{deleteOne} e Lote (\texttt{bulkWrite})}\par
      \begin{lstlisting}[language=mongosh]
// Remover 1 documento especifico:
db.books.deleteOne({ "isbn": "201" })
// Retorno: { acknowledged: true, deletedCount: 1 }

// Operacao em Lote (Batch):
db.books.bulkWrite([
  { insertOne: { document: {
      isbn: "202", name: "MongoDB Guide", price: 40
  }}},
  { updateOne: {
      filter: { isbn: "101" },
      update: { "$set": { price: 49 } }
  }}
])
      \end{lstlisting}
    \end{column}
  \end{columns}
\end{frame}
